\begin{center}
Problema H

Rosa dos Ventos

{\tiny Por Heverton Barros de Macêdo}

Timelimit: $1$	
\end{center}

Em uma competição de parapente há um série de provas, sendo a prova de navegação a que demanda maior estratégia por parte da equipe. Na prova de navegação, alguns minutos antes de seu início, os pilotos recebem um roteiro com os pontos a serem sobrevoados. Ganha a prova, o piloto que conseguir passar sobre a maior quantidade de pontos dentro do limite de tempo previamente estabelecido pela direção de prova.

Nesse ano a organização do Campeonato Goiano de Parapente (CGP) decidiu inovar exigindo que todos os pilotos façam somente curvas de 90 graus. Além disso, ao término da prova, antes de pousar, o piloto deve estar voando na mesma direção em que saiu da rampa.

O problema agora é que a direção de prova está com dificuldades para analisar os relatórios que contém as curvas realizadas pelos pilotos.

Ajude a CGP a avaliar o desempenho dos pilotos. Sua missão é construir um programa que possa analisar a sequência de curvas realizadas por um piloto e apresentar a direção em que ele terminou a prova.

\vspace{0.3cm}
\textbf{Entrada}

A entrada contém vários casos de teste. A primeira linha de cada caso corresponde a quantidade de comandos $(Q)$ realizada pelo piloto $(1 \leq Q \leq 100)$, assim como a indicação da posição $P$ (Norte, Sul, Leste ou Oeste) da rampa em que ele decolou $(P \in \left\{N, S, L, O\right\})$, sendo essas duas informações separadas por um único espaço. A segunda linha indica a sequência de comandos $(C)$ realizada pelo piloto, podendo ser $D$, $E$ ou $F$, com significado de curva à Direita, curva à Esquerda e siga em Frente, respectivamente $(C \in \left\{D, E, F\right\})$.

\vspace{0.3cm}
\textbf{Saída}

O programa deve produzir como saída uma única letra indicando a direção em que o piloto finalizou a prova de navegação, podendo ser $N$ (Norte), $S$ (Sul), $L$ (Leste) ou $O$ (Oeste).

O final da entrada é indicado por $Q = 0$.

\begin{table}[ht]
    \centering
    \begin{tabular}{|p{7cm}|p{7cm}|}
        \hline
        \begin{center}
            \vspace{-0.3cm}
            \textbf{Exemplos de Entrada}
            \vspace{-0.3cm}
        \end{center} 
         &  
        \begin{center}
            \vspace{-0.3cm}
            \textbf{Exemplos de Saída}
            \vspace{-0.3cm}
        \end{center} \\ \hline
         \texttt{1 N} & \texttt{O} \\  
         \texttt{E} & \texttt{L} \\
         \texttt{3 O} & \texttt{S} \\  
         \texttt{DDF} & \texttt{S} \\
         \texttt{7 S} & \texttt{N} \\  
         \texttt{DDFDEDD} & \texttt{} \\
         \texttt{7 N} & \texttt{} \\  
         \texttt{DFDFEDF} & \texttt{} \\
         \texttt{13 O} & \texttt{} \\  
         \texttt{FEFEFEFEFEFDD} & \texttt{} \\
         \texttt{0} &  \\ \hline
    \end{tabular}
    \caption{Exemplos de entradas e saídas}
    \label{tab:inputL}

    
\end{table}

\newpage

\begin{center}
\noindent
\lstinputlisting{abobora_22/H/H2.cpp}
\end{center}
